\section{部分子分布}

\subsection{一般矩阵元素与 Lorentz 不变性}

历史上 \cite{Feynman:1973xc}，PDFs 的引入
是为了描述 \mbox{自旋-1/2 夸克。}
由于与自旋相关的复杂性
并不影响
部分子分布这一概念本身，
我们从一个简单的标量理论例子
出发。在那种情形下，关于靶的信息
累积在一般矩阵元素 $\langle p | \phi (0) \phi(z) | p \rangle $ 中。
由 Lorentz 不变性，它是两个标量的函数：$(pz) \equiv - \nu$ 与 $z^2$
（若希望对类空 $z$ 取正值，则用 $-z^2$）：
\begin{align}
\langle p |   \phi(0) \phi (z)|p \rangle
=  & {\cal M} (-(pz), -z^2)
\, .
\label{lorentz}
\end{align}
可以证明 \cite{Radyushkin:2016hsy,Radyushkin:1983wh}，对所有有贡献的 Feynman 图，
它对 $(pz)$ 的 Fourier 变换 ${\cal P} (x, -z^2)$
具有 $-1 \leq x \leq 1$ 的支撑，即
\begin{align}
{\cal M} (-(pz), -z^2)
&   =
\int_{-1}^1 dx
\, e^{-i x (pz) } \,  {\cal P} (x, -z^2)  \   .
\label{MPD}
\end{align}
注意，式 (\ref{MPD}) 给出了 $x$ 的一个协变定义。
无需假设
$p^2=0$ 或 $z^2=0$ 即可定义 $x$。

\subsection{共线 PDFs}

选取 $p$ 和 $z$ 的一些特殊情形，
便可用同一个函数 ${\cal M} (-(pz), -z^2)$
表示各种部分子分布的表达式。
特别地，取类光 $z$，例如只具有光前减分量 $z_-$ 的情形，
我们用扭度-2 部分子分布 $f(x)$ 参数化该矩阵元素：
\begin{align}
{\cal M} (-p_+ z_- , 0)  =
\int_{-1}^1 dx \, f(x) \,
e^{-ixp_+ z_-} \,  \   ,
\label{twist2par0}
\end{align}
其中 $f(x)$ 具有通常的解释：它是部分子携带靶动量分量 $p_+$ 的份额为 $x$
的概率。
逆关系为
\begin{align}
f  (x)  =\frac{1}{2 \pi}  \int_{-\infty}^{\infty}  d\nu \, e^{-i x \nu}  \, {\cal M}(\nu , 0)    = {\cal P} (x, 0)
\  .
\label{fxMnu}
\end{align}
由于 $f(x) =   {\cal P} (x, 0) $，函数 $ {\cal P} (x, -z^2) $
把 PDFs 推广到非类光的间隔 $z^2$ 上，我们将称之为
{\it 伪部分子分布（pseudo-PDF）}。变量 $(pz)$ 称为
{\it Ioffe 时间} \cite{Ioffe:1969kf}，
而 ${\cal M}(\nu , -z^2) $ 是 {\it Ioffe 时间分布} \cite{Braun:1994jq}。

注意，$ {\cal P} (x, -z^2) $ 的定义比 $f(x)$ 的定义更简单，
因为它不需要取微妙的 $z^2 \to 0$ 极限。
在可重正理论中，
函数 $ {\cal M} (\nu, z^2)  $ 具有
$\sim \ln z^2$ 奇异性，
它们产生部分子密度的微扰演化。
在算符乘积展开
（OPE）框架内，
这些 $\ln z^2$ 奇异性通过某种方案（例如维数重正化）被减除，
所得 PDFs 依赖于重正化标度 $\mu$，即
$f(x) \to   f(x, \mu^2)$。

\subsection{横动量依赖分布}

当 $z^2$ 类空时，可以把 $-z^2$
视为二维矢量 $ \{z_1,z_2\}$ 的模方，
并对其分量引入二维 Fourier 变换，即写成
\begin{align}
{\cal P} (x,  z_1^2+z_2^2)
&   =
\int_{-\infty}^\infty  {d k_1 } e^{i k_1 z_1}
\nn & \times
\int_{-\infty}^\infty  dk_2 \,e^{ik_2z_2}
{\cal F} (x, k_1^2+k_2^2)
\   .
\label{PTMD}
\end{align}
由于 ${\cal P} (x,  z_1^2+z_2^2) $ 在 $ \{z_1,z_2\}$ 平面内
具有转动不变性，
函数 ${\cal F} (x, k_1^2+k_2^2)$ 只通过 $k_1^2+k_2^2$ 依赖
${k_1,k_2}$，这一点已经体现在记号中。
将此表示与式 (\ref{MPD}) 结合，得到
\begin{align}
{\cal M} (\nu,  z_1^2+z_2^2)
&   =  \int_{-1}^1 dx
\, e^{i x \nu  } \,
\int_{-\infty}^\infty  {d k_1 } e^{i k_1 z_1}
\nn & \times
\int_{-\infty}^\infty  dk_2 \,e^{ik_2z_2}
{\cal F} (x, k_1^2+k_2^2)
\   .
\label{MTMD}
\end{align}

${\cal F} (x, k_1^2+k_2^2)$ 的物理解释可以在靶动量 $p$ 为纵向的参考系中给出：
\mbox{$p= (E, {\bf 0}_\perp, P)$,} 而矢量
$\{z_1,z_2\} $ 位于横平面内。取只含
$z_-$ 和 $z_\perp$ 分量的 $z$，
便可以把 ${\cal F} (x, k_\perp^2)$ 认同于
{\it TMD} 并写成
\begin{align}
{\cal P} (x,  z_\perp^2)
&   =
\int  d^2{\bf k}_\perp   e^{i ({\bf k}_\perp {\bf z}_\perp)}
{\cal F} (x, {k}_\perp^2)
\   .
\label{MTMD0}
\end{align}
此时 pseudo-PDFs $ {\cal P} (x,  z_\perp^2) $ 与
{\it 碰撞参数分布（impact parameter distributions）} 相重合，
后者是 TMD 研究
中广泛使用的熟知概念。

${\cal M} (\nu,  z_\perp^2) $ 中的
\mbox{$\sim \ln z_\perp^2$} 项
来自 ${\cal F} (x, k_\perp^2)$ 的 $\sim 1/k_\perp^2$ 硬尾部。
因此，把
${\cal M} (\nu, z_\perp^2)  $ 视为软部分 ${\cal M}^{\rm soft} (\nu, z_\perp^2)$
（它具有有限的 $z_\perp^2\to 0$ 极限）
与反映演化的硬部分之和是有意义的。
对于 TMDs，软部分衰减得快于 $1/k_\perp^2$，比如像
Gaussian 型 $e^{-k_\perp^2/\Lambda^2}$。于是分布在 $z_\perp$ 空间中
集中在 \mbox{$z_\perp \sim 1/\Lambda$} 区域。
